\section{Motivation and statement}

A basic problem in bilipschitz invariant theory is the following.  Given a
Hilbert space $V$ and a subgroup $G\le O(V)$, when does the orbit space $V/G$,
with its natural quotient metric, admit a bilipschitz embedding into a Hilbert
space?  Cahill, Iverson, and Mixon formulated this program in
\cite{CahillIversonMixon}; their motivation comes from stable invariant
representations of data observed only modulo symmetries, such as relabelings of
graphs, permutations of point clouds, rotations of landmarks, and time shifts of
signals.  In that paper, Problem~48 asks specifically whether
$\ell^2(\Z)/\Z$ admits a bilipschitz embedding into Hilbert space, and asks for
the exact value of its Euclidean distortion \cite[Problem~48]{CahillIversonMixon}.

Throughout, all Hilbert spaces are real.  We focus on the sharper spherical
obstruction for the bilateral shift: the orbit space of the unit sphere already
has infinite Euclidean distortion.  No separate full-space statement is needed
below.

Let
\[
        \Sigma=\Sphere(\ell^2(\Z))/\langle S\rangle.
\]
The shift orbits on the unit sphere are closed.  Indeed, if $S^{k_j}x\to y$
in norm, then either $(k_j)$ has a constant subsequence, in which case $y$
belongs to the orbit of $x$, or $|k_j|\to\infty$ along a subsequence.  In the
latter case $S^{k_j}x\rightharpoonup0$, so norm convergence would force
$y=0$, impossible on the unit sphere.  Thus the following quotient
pseudometrics are genuine metrics.  The chordal quotient metric is
\[
        d_{\rm ch}([x],[y])=\inf_{k\in\Z}\norm{x-S^k y}_2.
\]
The intrinsic angular quotient metric is
\[
        d_{\rm ang}([x],[y])=\inf_{k\in\Z}\arccos\ip{x}{S^k y}.
\]
For unit vectors, $\norm{x-y}=2\sin(\theta(x,y)/2)$, where
$\theta(x,y)=\arccos\ip{x}{y}$.  Since $t\mapsto2\sin(t/2)$ is increasing on
$[0,\pi]$,
\begin{equation}\label{eq:metric-equivalence}
        d_{\rm ch}=2\sin\frac{d_{\rm ang}}2,
        \qquad
        \frac2\pi d_{\rm ang}\le d_{\rm ch}\le d_{\rm ang}.
\end{equation}
Thus either metric may be used for questions of Hilbert distortion.

\begin{theorem}\label{thm:main}
For the bilateral shift on $H=\ell^2(\Z)$,
\[
        c_2\bigl(\Sphere(H)/\langle S\rangle,d_{\rm ch}\bigr)=\infty.
\]
Consequently,
\[
        c_2\bigl(\Sphere(H)/\langle S\rangle,d_{\rm ang}\bigr)=\infty.
\]
Equivalently, the shift quotient of the Hilbert sphere contains finite subsets
of arbitrarily large Hilbert distortion.
\end{theorem}